\chapter{Firmware Architecture}

\section{Architectural Objectives}

The firmware architecture is designed to enforce structural clarity, execution determinism, and controlled modular growth.

The architectural objectives are:

\begin{itemize}
	\item Enforce separation between communication, control logic, and hardware interaction.
	\item Maintain deterministic execution ordering.
	\item Constrain inter-module dependencies.
	\item Enable incremental feature expansion without structural degradation.
	\item Isolate hardware-specific implementation details.
\end{itemize}

These objectives define the structural organization of the firmware.

% ------------------------------------------------

\section{Architectural Decomposition Model}

The firmware follows a layered decomposition strategy. Each layer encapsulates a specific abstraction level and interacts only with adjacent layers.

\subsection{Logical Layers}

The system is partitioned into four logical layers:

\begin{enumerate}
	\item Application / Control Layer
	\item Communication Layer
	\item Hardware Abstraction Layer (HAL)
	\item Physical Interface Layer
\end{enumerate}

\subsubsection*{Application / Control Layer}

This layer contains system-level logic and internal state management.  
It does not access hardware registers or communication buffers directly.

\subsubsection*{Communication Layer}

This layer manages structured packet exchange and data validation.  
It translates external serialized data into internal control variables and formats telemetry structures for transmission.

\subsubsection*{Hardware Abstraction Layer}

The HAL provides a stable interface between logical control variables and microcontroller peripherals.  
All direct peripheral configuration and register-level interaction is confined to this layer.

\subsubsection*{Physical Interface Layer}

Represents electrical connections and external components.  
No firmware logic exists at this layer.

% ------------------------------------------------

\section{Module-Level Decomposition}

Within the layered structure, the firmware is divided into discrete modules.  
Each module has a well-defined responsibility and limited external visibility.

Primary module categories include:

\begin{itemize}
	\item Communication Parsing Module
	\item Control Variable Management Module
	\item Motor Control Module
	\item Servo Control Module
	\item Telemetry Aggregation Module
	\item Sensor Interface Modules
	\item GPIO Expansion Driver
	\item Output Management Module
	\item Debug and Diagnostic Module
\end{itemize}

Modules interact through controlled data interfaces rather than direct cross-calls.

% ------------------------------------------------

\section{Execution Model}

The firmware executes within a single main control loop.  
No preemptive scheduler or real-time operating system is employed.

\subsection{Deterministic Ordering}

Each loop iteration executes subsystems in a fixed order.  
The ordering is intentionally static to guarantee predictable behavior and bounded execution time.

\subsection{Polling Strategy}

Subsystems are serviced through structured polling.  
This eliminates concurrency conflicts and reduces synchronization complexity.

\subsection{Non-Blocking Constraint}

Blocking delays are avoided within loop-executed modules.  
All time-dependent behavior must be implemented using state tracking rather than delay-based waiting.

This constraint ensures bounded worst-case loop duration.

% ------------------------------------------------

\section{Data Ownership and State Management}

The firmware maintains a centralized internal state model.

\subsection{State Variables}

State variables represent:

\begin{itemize}
	\item Target actuator values
	\item Measured sensor data
	\item Communication status
	\item Output states
\end{itemize}

Each module is permitted to modify only the state variables within its defined responsibility domain.

\subsection{State Update Discipline}

State updates occur in deterministic phases of the control loop.  
Read-modify-write patterns are controlled to avoid unintended side effects.

% ------------------------------------------------

\section{Data Flow Structure}

\subsection{Command Processing Path}

Incoming serialized data is:

\begin{enumerate}
	\item Received into a communication buffer
	\item Validated and parsed
	\item Translated into internal control variables
	\item Applied to actuator update logic
\end{enumerate}

\subsection{Telemetry Processing Path}

Telemetry data is:

\begin{enumerate}
	\item Collected from state variables and sensor modules
	\item Assembled into structured packets
	\item Serialized into a transmission buffer
	\item Sent through the communication interface
\end{enumerate}

Buffers isolate communication timing from control execution timing.

% ------------------------------------------------

\section{Interface Boundaries}

Strict interface boundaries are enforced between:

\begin{itemize}
	\item Communication and Control logic
	\item Control logic and HAL
	\item HAL and hardware peripherals
\end{itemize}

Direct cross-layer access is prohibited.  
This containment prevents structural coupling and simplifies future refactoring.

% ------------------------------------------------

\section{Hardware Abstraction Strategy}

All hardware-specific configuration is localized within the HAL.

Abstraction mechanisms include:

\begin{itemize}
	\item Encapsulated peripheral initialization
	\item Driver-level isolation
	\item Pin mapping centralization
	\item Replaceable hardware drivers
\end{itemize}

Control modules remain independent of specific peripheral registers or microcontroller configuration details.

% ------------------------------------------------

\section{Scalability and Controlled Expansion}

The architecture supports feature growth through modular extension.

\subsection{Sensor Integration}

New sensors are introduced via independent interface modules connected through the HAL.

\subsection{Actuator Integration}

New actuators are implemented as discrete control modules without altering communication parsing logic.

\subsection{I/O Expansion}

External expansion devices increase hardware capacity while maintaining logical separation from application logic.

\subsection{Structural Stability}

Architectural integrity is preserved by:

\begin{itemize}
	\item Maintaining layer boundaries
	\item Avoiding circular dependencies
	\item Restricting module visibility
\end{itemize}

These constraints prevent architectural degradation as the firmware evolves.